\documentclass{report}
\usepackage[utf8]{inputenc}

\title{Practica 1 \\ Introducción a las herramientas del laboratorio}
\author{iabin Arteaga}
\date{August 2019}

\usepackage{natbib}
\usepackage{graphicx}
\usepackage{hyperref}
\usepackage{longtable}
\usepackage{booktabs}

\begin{document}

{Práctica 2}

{Tipos básicos, números y cadenas }

{}

{Objetivo}

{El objetivo de esta práctica es que conozcas los tipos básicos del lenguaje java, que aprendas a utilizar los distintos tipos numéricos del lenguaje, aprender a diferenciar entre un tipo básico y un objeto y aprender que es una cadena y cómo operar con ellas, utilizando las funciones de su clase.}

{}

{Desarrollo}

{Descarga la clase IntAndString.java }

{La práctica consiste en completar los siguientes métodos:}

{}

\href{}{}\href{}{}

\begin{longtable}[c]{@{}ll@{}}
\toprule
\begin{minipage}[t]{0.47\columnwidth}\raggedright\strut
{FUNCIÓN}
\strut\end{minipage} & \begin{minipage}[t]{0.47\columnwidth}\raggedright\strut
{EXPLICACION}
\strut\end{minipage}\tabularnewline
\begin{minipage}[t]{0.47\columnwidth}\raggedright\strut
{producto}
\strut\end{minipage} & \begin{minipage}[t]{0.47\columnwidth}\raggedright\strut
{Producto de dos números}
\strut\end{minipage}\tabularnewline
\begin{minipage}[t]{0.47\columnwidth}\raggedright\strut
{division}
\strut\end{minipage} & \begin{minipage}[t]{0.47\columnwidth}\raggedright\strut
{División de dos números}
\strut\end{minipage}\tabularnewline
\begin{minipage}[t]{0.47\columnwidth}\raggedright\strut
{suma}
\strut\end{minipage} & \begin{minipage}[t]{0.47\columnwidth}\raggedright\strut
{Suma de dos números}
\strut\end{minipage}\tabularnewline
\begin{minipage}[t]{0.47\columnwidth}\raggedright\strut
{potencia}
\strut\end{minipage} & \begin{minipage}[t]{0.47\columnwidth}\raggedright\strut
{Potencia de números }
\strut\end{minipage}\tabularnewline
\begin{minipage}[t]{0.47\columnwidth}\raggedright\strut
{raizCuadrada}
\strut\end{minipage} & \begin{minipage}[t]{0.47\columnwidth}\raggedright\strut
{Raíz cuadrada de un número}
\strut\end{minipage}\tabularnewline
\begin{minipage}[t]{0.47\columnwidth}\raggedright\strut
{chicharronera }
\strut\end{minipage} & \begin{minipage}[t]{0.47\columnwidth}\raggedright\strut
{Solo el resultado de la parte positiva de la raiz}
\strut\end{minipage}\tabularnewline
\begin{minipage}[t]{0.47\columnwidth}\raggedright\strut
{hipotenusa}
\strut\end{minipage} & \begin{minipage}[t]{0.47\columnwidth}\raggedright\strut
{Recibe 2 lados, regresa la hipotenusa}
\strut\end{minipage}\tabularnewline
\begin{minipage}[t]{0.47\columnwidth}\raggedright\strut
{concatenacion }
\strut\end{minipage} & \begin{minipage}[t]{0.47\columnwidth}\raggedright\strut
{Une dos cadenas}
\strut\end{minipage}\tabularnewline
\begin{minipage}[t]{0.47\columnwidth}\raggedright\strut
{ultimaLetra}
\strut\end{minipage} & \begin{minipage}[t]{0.47\columnwidth}\raggedright\strut
{Devuelve la última letra}
\strut\end{minipage}\tabularnewline
\begin{minipage}[t]{0.47\columnwidth}\raggedright\strut
{longitud}
\strut\end{minipage} & \begin{minipage}[t]{0.47\columnwidth}\raggedright\strut
{Devuelve la longitud de la cadena}
\strut\end{minipage}\tabularnewline
\begin{minipage}[t]{0.47\columnwidth}\raggedright\strut
{reemplazaLasA}
\strut\end{minipage} & \begin{minipage}[t]{0.47\columnwidth}\raggedright\strut
{Reemplaza todas las A por otra cadena}
\strut\end{minipage}\tabularnewline
\begin{minipage}[t]{0.47\columnwidth}\raggedright\strut
{parte}
\strut\end{minipage} & \begin{minipage}[t]{0.47\columnwidth}\raggedright\strut
{Parte una cadena dado un separador}
\strut\end{minipage}\tabularnewline
\begin{minipage}[t]{0.47\columnwidth}\raggedright\strut
{quitaEspacios}
\strut\end{minipage} & \begin{minipage}[t]{0.47\columnwidth}\raggedright\strut
{Quita los espacios finales e iniciales a una cadena}
\strut\end{minipage}\tabularnewline
\bottomrule
\end{longtable}

{}

{La práctica tiene un conjunto de pruebas, de tal manera que cuando ejecutes y compiles el programa, la terminal te dirá qué métodos están mal.}

\IfFileExists{images/image1.png}{\includegraphics{images/image1.png}}{}

{Cuando tus métodos pasen las pruebas, y vuelvas a ejecutar el programa, igual la terminal te lo dirá.}

\IfFileExists{images/image2.png}{\includegraphics{images/image2.png}}{}

{}

{El formato de entrega es exactamente igual que el anterior}

{Crean una carpeta que se llame practica02 y la colocan en su carpeta de trabajo de ICC}

{Crean una carpeta que se llame src y ahí colocan su código, anexar en el README de la práctica instrucciones de como compilar y ejecutar y luego simplemente la suben a gihub, }

{Esto será así para todas las siguientes prácticas y proyectos.}

{El repositorio se vería similar a esto.}

\begin{verbatim}
./
|-- practica01
|-- practica02
|   |-- README.txt
|   `-- IntAndString.java
`-- README.txt
\end{verbatim}

{}

\end{document}
